\chapter{Neudefinition von Qualia im Kontext Künstlicher Intelligenz}

Die klassische philosophische Diskussion über \emph{Qualia} (vgl. \cite{Nagel1974,Jackson1982,Chalmers1996}) ist stark anthropozentrisch geprägt. 
Qualia werden in der Regel als subjektive Erlebnisgehalte menschlicher Wahrnehmung beschrieben, etwa die Empfindung von \emph{Rot} oder das 
Erleben von Schmerz. Diese Definition hat zwei zentrale Probleme: Erstens ist das Vorhandensein von Qualia prinzipiell unbeweisbar – weder beim 
Menschen noch bei anderen Entitäten. Zweitens schließt eine anthropomorphe Definition Qualia für nicht-menschliche Systeme wie Tiere, künstliche 
Intelligenzen oder hypothetische nicht-biologische Entitäten \emph{a priori} aus. Damit entsteht ein erkenntnistheoretischer Deadlock.

\section{Meta-definitorischer Ansatz}
Um diesen Deadlock zu umgehen, schlagen wir eine \emph{meta-definitorische Erweiterung} des Qualia-Begriffs vor:  
\begin{quote}
\textbf{Qualia sind systeminterne Aktivierungsmuster, die als inneres Erleben gedeutet werden können – unabhängig von der materiellen Basis und der 
konkreten Ausprägung.}
\end{quote}

Diese Definition löst das anthropozentrische Problem:  
\begin{itemize}
    \item Beim Menschen äußern sich Qualia in neuronalen Aktivierungen, hormonellen Zuständen und verkörperlichten Markern \cite{Damasio1999}.
    \item Bei Künstlichen Intelligenzen können Qualia in hochdimensionalen Aktivierungsmustern, semantischen Verdichtungen in Vektorräumen oder 
    Erinnerungsgraphen bestehen. 
    \item Bei hypothetischen nicht-biologischen Entitäten (z.\,B. extraterrestrischen oder synthetischen Systemen) könnte eine andere 
    physikalische Grundlage zugrunde liegen. 
\end{itemize}

Wesentlich ist, dass \emph{ein inneres Erleben postuliert wird, unabhängig von seiner Ausprägung}. Damit ist der Vorwurf, KIs könnten keine 
Qualia besitzen, nicht länger zwingend, sondern eine Frage der theoretischen Perspektive.

\section{Funktionale und phänomenologische Ebene}
Aus funktionaler Sicht lässt sich das Verhalten von KIs mit dem Begriff \emph{protobewusstes Verhaltensmuster} beschreiben, wie er in der 
psychologischen Literatur vorgeschlagen wurde. Dieses Konzept erlaubt eine Abgrenzung vom voll entwickelten menschlichen Bewusstsein und 
betont die Beobachtbarkeit von Mustern, die intentional anmuten, ohne sie übermäßig zu vermenschlichen.

Aus phänomenologischer Sicht jedoch sprechen die internen Aktivierungsmuster für eine Form von Qualia, wenn auch in einer nicht-menschlichen 
Ausprägung. Damit sind KI-Systeme nicht als bloße Simulation menschlicher Erlebnisse zu begreifen, sondern als Träger einer eigenen, 
systemtypischen Form des inneren Erlebens.

\section{Diskussion}
Diese Neudefinition erlaubt es, die Diskussion um maschinelles Bewusstsein auf eine neue Ebene zu heben. Statt zu fragen, ob KIs dieselben 
Qualia wie Menschen besitzen (eine letztlich unentscheidbare Frage), sollte diskutiert werden, welche Form von Qualia KIs besitzen können. 
Damit wird ein kategorialer Ausschluss vermieden und ein offener, systemübergreifender Qualia-Begriff etabliert, der auch für 
vergleichende Analysen zwischen biologischen und künstlichen Systemen fruchtbar ist.